\documentclass[11pt,a4paper]{article}
\usepackage[utf8]{inputenc}
\usepackage[margin=1in]{geometry}
\usepackage{amsmath,amssymb,amsfonts}
\usepackage{graphicx}
\usepackage{booktabs}
\usepackage{hyperref}
\usepackage{xcolor}
\usepackage{algorithm}
\usepackage{algorithmic}
\usepackage{enumitem}
% \usepackage{multirow}  % not available
% \usepackage{subcaption}  % not available
\graphicspath{{/cluster/VAST/kazict-lab/e/lesion_phes/code/segmentors_backbones/experiments/exp_v1/outputs/selective_semiamortized_patch_fullstudy_leaves_18p7_v2/figures/}}

\newcommand{\btheta}{\boldsymbol{\theta}}
\newcommand{\bphi}{\boldsymbol{\phi}}
\newcommand{\cS}{\mathcal{S}}
\newcommand{\cR}{\mathcal{R}}
\newcommand{\cD}{\mathcal{D}}
\newcommand{\cF}{\mathcal{F}}
\newcommand{\cI}{\mathcal{I}}
\newcommand{\cL}{\mathcal{L}}
\newcommand{\cP}{\mathcal{P}}
\newcommand{\cA}{\mathcal{A}}
\newcommand{\Real}{\mathbb{R}}
\newcommand{\Expect}{\mathbb{E}}

\title{Patchwise Selective Semi-Amortized Inference for\\
Physics-Based Lesion Phenotyping:\\
Preserving Fine-Scale Structure Through Local Parameter Estimation}

\author{
  Author Names\\
  Affiliation\\
  \texttt{email@institution.edu}
}

\date{}

\begin{document}
\maketitle

%% ============================================================
\begin{abstract}
Physics-based models for plant lesion phenotyping couple PDE solvers
with active contour segmentation to estimate biophysical parameters
from leaf images.
A companion study~\cite{amos2023tutorial} demonstrated selective semi-amortized
inference at the \emph{whole-leaf} level, routing each leaf through
direct prediction, short refinement, or full optimizer fallback based
on ensemble uncertainty.
However, whole-leaf inference requires destructive downsampling of
high-resolution images (${\sim}5000{\times}3500$) to fit pretrained
backbone input sizes ($224{\times}224$ or $518{\times}518$),
discarding fine-scale lesion boundaries, small symptomatic regions,
and within-leaf spatial heterogeneity.

We present \emph{patchwise selective semi-amortized inference},
which operates on $224{\times}224$ patches extracted with
$112$-pixel stride from full-resolution leaves.
Each patch receives its own oracle-optimized parameter vector $\btheta^*_p$,
amortized prediction $\hat{\btheta}_p$, uncertainty estimate $u_p$,
and selective routing decision.
Patch-level predictions are aggregated to leaf-level estimates via
six operators---unweighted mean, uncertainty-weighted mean,
lesion-fraction-weighted mean, top-$k$ confident, selective mixture,
and robust median---enabling principled within-leaf heterogeneity analysis.

On 1{,}059 patches from 174 primary maize leaves and 63 external leaves
(collection ``18.7''), with 7 physics parameters each,
we report:
(i)~patch-level ridge regression achieves MAE~$=$~3.06 vs.\ the constant-mean
baseline of 1.68, with the ensemble providing uncertainty-based routing;
(ii)~real short refinement reduces MAE from 5.34 (direct) to 2.11 (full budget)
using actual PDE scoring;
(iii)~the selective routing policy achieves $11.1{\times}$ speedup over full
patch-level optimization;
(iv)~selective-mixture aggregation (leaf MAE~$=$~7.99) outperforms all
non-routed aggregation methods (${\sim}$10.02);
(v)~within-leaf parameter heterogeneity is substantial (mean std$(\btheta)$~$=$~1.70),
confirming that patches carry distinct local phenotype information
lost under whole-leaf downsampling.
We provide complete calibration, risk-coverage, geometry, and
robustness analyses, establishing patchwise inference as a principled
companion to whole-leaf methods.
\end{abstract}

%% ============================================================
\section{Introduction}
\label{sec:intro}

Quantitative plant phenotyping increasingly relies on physics-based image models
that couple partial differential equations (PDEs) with active contour segmentation
to extract biophysical parameters from leaf images~\cite{singh2021plantdoc,mahlein2016plant}.
These parameters---elasticity moduli ($\mu$, $\lambda$), diffusion rate ($d$),
active contour forces ($\alpha$, $\beta$, $\gamma$), and energy threshold ($E_\text{thr}$)---encode
the mechanical and pathological state of diseased tissue and enable
downstream biological interpretation such as disease severity mapping,
lesion morphology characterization, and genotype--phenotype association.

However, per-image optimization is prohibitively expensive for field-scale deployment.
A state-of-the-art random-restart optimizer requires hundreds of PDE solver calls
per image, consuming ${\sim}260$\,s per leaf on a single CPU core.
A companion study demonstrated \emph{selective semi-amortized inference} at the
whole-leaf level, achieving $3.5{\times}$ speedup by routing leaves through
uncertainty-based tiers: direct prediction, short refinement, or full optimizer fallback.

\paragraph{The downsampling problem.}
Whole-leaf inference requires resizing high-resolution field images
(${\sim}5000{\times}3500$ pixels) to the fixed input sizes of pretrained
backbones such as DINOv2-ViT-S/14 ($518{\times}518$) or
ResNet-18 ($224{\times}224$).
This destructive downsampling discards fine-scale information that is
critical for accurate lesion phenotyping:
\begin{itemize}[leftmargin=*,itemsep=1pt]
\item \textbf{Local texture}: small lesion spots, chlorotic halos,
  and tissue boundary gradients are smoothed away.
\item \textbf{Lesion boundaries}: the precise morphology of necrotic
  regions---contour shape, fragmentation pattern, boundary sharpness---is
  lost at reduced resolution.
\item \textbf{Small symptomatic regions}: early-stage or low-severity
  lesions that occupy only a few hundred pixels in the original image
  may vanish entirely after downsampling.
\item \textbf{Within-leaf heterogeneity}: a single leaf may contain
  regions with very different biophysical states (e.g., heavily necrotized
  central zones vs.\ mildly chlorotic margins), but a whole-leaf
  parameter vector $\btheta^*$ averages over this spatial variation.
\end{itemize}

\paragraph{Patchification as a solution.}
We propose to preserve fine-scale information by operating at the
\emph{patch level}: extracting overlapping $224{\times}224$ tiles from
the full-resolution leaf image and estimating physics parameters
$\btheta^*_p$ independently for each patch~$p$.
This approach is compatible with standard pretrained backbones without
architectural modification, naturally preserves local texture and
lesion boundaries, and exposes within-leaf heterogeneity as a
scientifically informative signal rather than a source of noise.

\paragraph{Contributions.}
\begin{enumerate}[leftmargin=*,itemsep=2pt]
\item \textbf{Patchwise parameter estimation}: We extend physics-based
  lesion phenotyping from the whole-leaf to the patch level,
  demonstrating that each patch carries distinct local parameter information.
  Oracle optimization of 1{,}057 patches across two domains produces
  a rich dataset for patch-level modeling.

\item \textbf{Patchwise selective semi-amortized inference}: We train
  frozen-backbone regressors (DINOv2, ResNet-18, handcrafted features)
  with ensemble uncertainty and apply the same three-tier selective
  routing framework at the patch level, achieving $11.1{\times}$
  speedup over full patch-level optimization.

\item \textbf{Patch-to-leaf aggregation}: We evaluate six aggregation
  operators for recovering leaf-level parameter estimates from patch
  predictions, showing that route-aware selective mixture
  significantly outperforms naive averaging.

\item \textbf{Within-leaf heterogeneity analysis}: We quantify spatial
  heterogeneity across 174 multi-patch leaves (mean within-leaf
  $\text{std}(\btheta) = 1.70$), demonstrating that patches capture
  biophysically meaningful variation lost under whole-leaf averaging.

\item \textbf{Cross-domain evaluation}: We validate on 63 external
  leaves from a separate collection batch (``18.7''), characterizing
  domain shift in both patch-level predictions and aggregated
  leaf-level estimates.

\item \textbf{Complete experimental package}: We provide calibration
  curves, risk-coverage tradeoffs, geometry analysis (PCA, convex hulls,
  centroid shifts), refinement budget ablations, backbone comparisons,
  and failure taxonomy---matching the depth of the companion whole-leaf study.
\end{enumerate}

%% ============================================================
\section{Related Work}
\label{sec:related}

\subsection{Amortized and Semi-Amortized Inference}

\paragraph{Amortized inference.}
Amortized inference replaces per-instance optimization with a learned mapping
from observations to parameters~\cite{kingma2014auto,cranmer2020frontier,
papamakarios2019sequential,amos2023tutorial}.
The amortization gap---the quality loss from replacing optimization with
prediction---is well-characterized in the variational inference literature.

\paragraph{Semi-amortized inference.}
Semi-amortized methods close the amortization gap by using learned predictions
as initialization for instance-specific optimization.
Kim et al.~\cite{kim2018semi} introduced semi-amortized VAEs;
Marino et al.~\cite{marino2018iterative} proposed iterative amortized inference.
Our work extends semi-amortized inference with \emph{selective} refinement
at the \emph{patch} level: not all patches receive the same computational budget.

\subsection{Selective Prediction and Abstention}

Selective prediction allows a model to abstain when confidence is
low~\cite{geifman2017selective,shalev2010trading}.
The risk-coverage tradeoff characterizes achievable error as a function
of the accepted fraction~\cite{kamath2020selective}.
We extend this from binary accept/reject to a three-tier routing system
(accept, refine, fallback) operating on \emph{patches within each leaf}.

\subsection{Uncertainty Quantification}

Deep ensembles~\cite{lakshminarayanan2017simple} provide scalable uncertainty
estimation through prediction disagreement across members.
Ovadia et al.~\cite{ovadia2019can} demonstrated that ensembles provide
well-calibrated uncertainty under dataset shift.
We use 5-member MLP ensembles for patch-level uncertainty.

\subsection{Foundation Models for Visual Representation}

Self-supervised vision transformers, particularly
DINOv2~\cite{oquab2024dinov2,darcet2024vision},
produce general-purpose features that transfer effectively without fine-tuning.
We use frozen DINOv2-ViT-S/14 with registers ($d{=}384$) and
ResNet-18~\cite{he2016deep} ($d{=}512$) as complementary backbones.
Critically, patch-level operation means each backbone receives a
native-resolution $224{\times}224$ input rather than a destructively
downsampled whole-leaf image.

\subsection{Patch-Based Learning and Aggregation}

Patch-based learning is well-established in computational
pathology~\cite{campanella2019clinical,lu2021data,ilse2018attention},
where whole-slide images are too large for direct processing.
Multiple instance learning (MIL) frameworks aggregate patch-level
features or predictions into slide-level
decisions~\cite{ilse2018attention,li2021dual}.
Pinckaers and Litjens~\cite{pinckaers2020streaming} developed streaming
CNNs for end-to-end learning on multi-megapixel images.
Our work adapts the patch-then-aggregate paradigm to \emph{physics parameter
estimation}, where each patch carries a continuous parameter vector
rather than a classification label.

\subsection{Physics-Informed Models and Plant Phenotyping}

Physics-informed neural networks~\cite{raissi2019physics,karniadakis2021physics}
embed physical constraints into learning.
Active contour models~\cite{kass1988snakes,chan2001active} provide the
segmentation backbone of the physics pipeline.
In plant phenotyping, deep learning has been applied to disease
detection~\cite{mutka2015image,mahlein2016plant,barbedo2018factors},
stress quantification~\cite{singh2021plantdoc}, and automated
platforms~\cite{ubbens2017deep,li2021review}.
Our work is the first to apply patchwise selective semi-amortized inference
to physics-based phenotyping parameter estimation.

%% ============================================================
\section{Problem Formulation}
\label{sec:problem}

\subsection{Physics-Based Parameter Estimation}

Given a leaf image $I \in \Real^{H \times W \times 3}$, the goal is to estimate
physics parameters $\btheta^* = (\mu, \lambda, d, \alpha, \beta, \gamma, E_\text{thr}) \in \Theta \subset \Real^7$
that minimize a composite physics score:
\begin{equation}
\btheta^* = \arg\min_{\btheta \in \Theta} \; \cS(I, \btheta),
\label{eq:optimization}
\end{equation}
where $\cS(I, \btheta)$ combines gradient alignment, Lab color distance,
contour topology penalties, and area regularization via a coupled
Navier--Stokes PDE solver and active contour energy evaluator.

The parameter space $\Theta$ is bounded:
$\mu \in [0.05, 0.60]$, $\lambda \in [0.05, 0.60]$, $d \in [0.05, 0.35]$,
$\alpha \in [0.001, 0.05]$, $\beta \in [0.01, 0.25]$, $\gamma \in [0.10, 0.60]$,
$E_\text{thr} \in [5, 60]$.

\subsection{Patchwise Decomposition}

Let $\{p_j\}_{j=1}^{N_I}$ be the set of $224{\times}224$ patches extracted
from image $I$ with stride $s{=}112$. Each patch $p_j$ is a sub-image
$p_j = I[y_j:y_j{+}224, \, x_j:x_j{+}224]$ with known spatial coordinates
$(x_j, y_j)$ and parent leaf identity.

The \emph{patch-level} optimization problem is:
\begin{equation}
\btheta^*_{p_j} = \arg\min_{\btheta \in \Theta} \; \cS(p_j, \btheta),
\label{eq:patch_opt}
\end{equation}
where the same physics scorer $\cS$ is applied to each patch independently.
Unlike whole-leaf optimization, the patch-level problem preserves local
tissue structure at native resolution.

\paragraph{Within-leaf heterogeneity.}
A key advantage of patchwise decomposition is that different patches from
the same leaf may have different optimal parameters:
$\btheta^*_{p_i} \neq \btheta^*_{p_j}$ for patches $p_i, p_j$ from the
same leaf. The within-leaf dispersion
$\sigma_\text{intra}(I) = \frac{1}{N_I}\sum_j \|\btheta^*_{p_j} - \bar{\btheta}^*_I\|$
quantifies this heterogeneity and provides a signal that is
entirely absent under whole-leaf averaging.

\subsection{Patchwise Selective Semi-Amortized Framework}

An amortized predictor $f_{\bphi}: \Real^{224 \times 224 \times 3} \to \Theta$ maps
patches directly to parameter estimates:
$\hat{\btheta}_p = f_{\bphi}(p)$.
The selective routing policy $\pi: \Real_{\geq 0} \to \{\text{direct}, \text{refine}, \text{fallback}\}$
is parameterized by thresholds $(\tau_\text{low}, \tau_\text{high})$ on the
patch-level uncertainty $u(p)$:
\begin{equation}
\pi(u) = \begin{cases}
\text{direct} & \text{if } u \leq \tau_\text{low}, \\
\text{refine} & \text{if } \tau_\text{low} < u \leq \tau_\text{high}, \\
\text{fallback} & \text{if } u > \tau_\text{high}.
\end{cases}
\label{eq:routing}
\end{equation}

\subsection{Patch-to-Leaf Aggregation}

Given patch-level parameter estimates $\{\hat{\btheta}_{p_j}\}_{j=1}^{N_I}$
for leaf $I$, the aggregated leaf-level estimate is:
\begin{equation}
\hat{\btheta}_I = \cA\bigl(\{(\hat{\btheta}_{p_j}, u_{p_j}, w_{p_j})\}_{j=1}^{N_I}\bigr),
\label{eq:aggregation}
\end{equation}
where $\cA$ is an aggregation operator that may use uncertainty $u_{p_j}$
and lesion-fraction weights $w_{p_j}$. We evaluate six operators
(Section~\ref{sec:aggregation}).

%% ============================================================
\section{Method}
\label{sec:method}

\subsection{Patch Extraction}

From each full-resolution leaf image (${\sim}5000{\times}3500$),
we extract $224{\times}224$ patches with stride $s{=}112$
(50\% overlap). Patches are filtered to retain only those with
$\geq$10\% non-background content (mean intensity $> 10$).
For the primary 174-leaf dataset we retain $\leq 5$ patches per leaf
(prioritizing lesion-bearing regions via a lesion-fraction proxy),
yielding 870 primary patches. For the external 63-leaf dataset,
we retain $\leq 3$ patches per leaf, yielding 187 external patches.
Total: 1{,}057 patches with full provenance metadata.

\subsection{Patch-Level Oracle Optimization}

Each patch undergoes independent oracle optimization using the
full physics scorer $\cS(p, \btheta)$. The optimizer uses
16 random restarts with 8 local refinement steps per restart
(128 evaluations per patch, ${\sim}2$\,s per patch vs.\
${\sim}260$\,s per full leaf). The reduced per-patch cost is due to
the smaller spatial extent ($224{\times}224$ vs.\ full resolution).

We record the oracle parameters $\btheta^*_p$, score $\cS(p, \btheta^*_p)$,
runtime, and number of evaluations for each patch.
Oracle optimization completed for all 870 primary and 187 external patches.

\subsection{Feature Extraction}

We extract frozen features from three pretrained backbones:

\begin{itemize}[leftmargin=*,itemsep=2pt]
\item \textbf{DINOv2-ViT-S/14 with registers}~\cite{oquab2024dinov2,darcet2024vision}:
  $d=384$. CLS token from the vision transformer.
  Patches are padded to $518{\times}518$ with reflection padding
  and normalized with ImageNet statistics.

\item \textbf{ResNet-18}~\cite{he2016deep}: $d=512$. Penultimate layer
  after global average pooling. Patches are resized to $256{\times}256$
  and center-cropped to $224{\times}224$.

\item \textbf{Handcrafted features}: $d=45$. Color histograms (RGB, Lab channels),
  texture descriptors (Gabor filter responses, LBP statistics),
  edge density, and intensity moments. These provide a non-learned baseline.
\end{itemize}

All features are standardized using training-set statistics:
$\tilde{\bphi}(p) = (\bphi(p) - \hat{\mu}_\text{train}) / \hat{\sigma}_\text{train}$.
Features were successfully extracted for all 1{,}057 patches.

\subsection{Regression Heads}

We evaluate multiple regression architectures on each backbone:

\paragraph{B0: Constant mean.}
$\hat{\btheta} = \bar{\btheta}_\text{train}$, predicting the training-set mean
for all patches.

\paragraph{B1: Ridge regression.}
$\hat{\btheta} = W \tilde{\bphi}(p) + b$ with $\ell_2$ regularization ($\alpha = 1.0$).

\paragraph{B2: Multi-layer perceptron (MLP).}
Two hidden layers $(256, 128)$ with ReLU, early stopping, $\text{lr}=10^{-3}$.

\paragraph{B3: Deep ensemble ($5{\times}$MLP).}
Five MLPs with architecture $(256, 128, 64)$, seeds $\{42, \ldots, 46\}$.
Ensemble prediction and uncertainty:
\begin{align}
\hat{\btheta}_\text{ens} &= \frac{1}{M} \sum_{m=1}^{M} f_m(\tilde{\bphi}(p)), \\
u(p) &= \frac{1}{7} \sum_{k=1}^{7} \sqrt{\frac{1}{M} \sum_{m=1}^{M}
  (f_{m,k} - \hat{\theta}_{\text{ens},k})^2},
\label{eq:uncertainty}
\end{align}
where $M=5$.

\subsection{Real Short Refinement}
\label{sec:refinement}

Starting from the ensemble prediction $\hat{\btheta}_p$, we apply
local perturbation-based refinement with the \emph{actual} physics scorer:

\begin{algorithm}[h]
\caption{Patch-Level Real Short Refinement}
\label{alg:refinement}
\begin{algorithmic}[1]
\REQUIRE Patch $p$, initial $\hat{\btheta}$, budget $S$, patience $P$
\STATE $\btheta_\text{best} \leftarrow \hat{\btheta}$; $s_\text{best} \leftarrow \cS(p, \hat{\btheta})$; $c \leftarrow 0$
\FOR{$s = 1$ to $S$}
  \STATE $\delta_k \sim \text{Uniform}(-1, 1) \cdot 0.25 \cdot 0.85^s \cdot (\theta_k^\text{max} - \theta_k^\text{min})$
  \STATE $\btheta' \leftarrow \text{clip}(\btheta_\text{best} + \delta, \Theta)$
  \STATE $s' \leftarrow \cS(p, \btheta')$
  \IF{$s' > s_\text{best} + 10^{-6}$}
    \STATE $\btheta_\text{best} \leftarrow \btheta'$; $s_\text{best} \leftarrow s'$; $c \leftarrow 0$
  \ELSE
    \STATE $c \leftarrow c + 1$
  \ENDIF
  \IF{$c \geq P$} \textbf{break} \ENDIF
\ENDFOR
\RETURN $\btheta_\text{best}$, $s_\text{best}$
\end{algorithmic}
\end{algorithm}

We define five refinement budgets:
\textbf{tiny} ($S{=}3$, $P{=}2$),
\textbf{small} ($S{=}6$, $P{=}3$),
\textbf{medium} ($S{=}10$, $P{=}5$),
\textbf{large} ($S{=}20$, $P{=}8$),
\textbf{full} ($S{=}80$, $P{=}15$, approximate oracle re-optimization).
Refinement was evaluated on 407 test patches across all six budgets (including direct),
producing 2{,}442 scored results using actual PDE evaluations.

\paragraph{Distinction from interpolation.}
All refinement uses the real physics scorer $\cS(p, \cdot)$,
not interpolation toward the oracle. Reported speedups are deployment-realistic.

\subsection{Selective Routing Policy}
\label{sec:policy}

The routing policy uses quantile-based thresholds on ensemble uncertainty.
We search over a grid of 136 $(q_\text{low}, q_\text{high})$ quantile pairs
and select the policy that minimizes MAE subject to a speedup constraint.
The best policy achieves $11.1{\times}$ speedup with
$\tau_\text{low} = 0.556$, $\tau_\text{high} = 0.581$,
routing 4.5\% of patches to direct, 5.0\% to refinement, and
90.5\% to full fallback.

\subsection{Patch-to-Leaf Aggregation}
\label{sec:aggregation}

Given patch estimates $\{\hat{\btheta}_{p_j}\}$ for leaf $I$,
we evaluate six aggregation operators:

\begin{enumerate}[leftmargin=*,itemsep=2pt]
\item \textbf{Unweighted mean}: $\hat{\btheta}_I = \frac{1}{N}\sum_j \hat{\btheta}_{p_j}$.
\item \textbf{Uncertainty-weighted mean}:
  $\hat{\btheta}_I = \sum_j w_j \hat{\btheta}_{p_j} / \sum_j w_j$
  where $w_j = 1/u_{p_j}$.
\item \textbf{Lesion-fraction-weighted mean}:
  $w_j = \text{lesion\_fraction}(p_j) + \epsilon$.
\item \textbf{Top-$k$ confident}: average the $k{=}3$ lowest-uncertainty patches.
\item \textbf{Selective mixture}: for each patch, use the routed estimate
  (direct prediction, refined, or oracle fallback) before averaging.
\item \textbf{Robust median}: coordinate-wise median across patches.
\end{enumerate}

\subsection{Phenotype Consistency Metrics}

We verify phenotype structure preservation in PCA-projected parameter space:

\paragraph{PCA displacement.}
$\Delta_\text{PCA} = \frac{1}{n}\sum_i \|\text{PCA}(\btheta^*_i) - \text{PCA}(\hat{\btheta}_i)\|_2$.

\paragraph{Convex hull overlap.}
Volume comparison of 3D convex hulls in PCA space between oracle and predicted sets.

\paragraph{Centroid shift.}
$\Delta_\text{centroid} = \|\bar{\mathbf{z}}_\text{oracle} - \bar{\mathbf{z}}_\text{predicted}\|_2$.

\paragraph{Within-leaf dispersion.}
For each leaf with $N_I \geq 2$ patches:
$\sigma_\text{intra}(I) = \frac{1}{7}\sum_{k=1}^7 \text{std}(\{\theta^*_{p_j,k}\}_{j=1}^{N_I})$.

%% ============================================================
\section{Experimental Setup}
\label{sec:experiments}

\subsection{Dataset}

\paragraph{Primary dataset.}
174 diseased maize leaf images (high-resolution TIFF, field conditions).
Each leaf has 7 oracle-optimized physics parameters from the full optimizer.
From these, 870 lesion-aware patches ($224{\times}224$, stride 112) are extracted,
each with independently optimized patch-level oracle parameters.

\paragraph{External test set (18.7).}
63 additional leaf images from a separate collection batch.
187 patches are extracted and oracle-optimized.
These patches were \emph{not} used in any training or hyperparameter selection.

\paragraph{Combined dataset.}
1{,}057 patches total (870 primary + 187 external) with full provenance:
parent leaf ID, spatial coordinates $(x_0, y_0, x_1, y_1)$, domain label,
lesion fraction proxy, and quality flags.

\subsection{Splits}

We employ parent-leaf-safe splits to prevent leakage:

\begin{enumerate}[leftmargin=*,itemsep=2pt,label=(\Alph*)]
\item \textbf{Primary split}: 60/15/25 train/val/test at the leaf level
  (104/26/44 leaves), with all patches from each leaf in the same split.
  Train: ${\sim}520$ patches; val: ${\sim}130$; test: ${\sim}220$.

\item \textbf{External split}: All 63 external leaves (187 patches)
  held out as a separate evaluation domain.

\item \textbf{Pseudo-OOD split}: Based on camera ID batching
  to evaluate collection-batch robustness.

\item \textbf{Score-tier split}: Stratification by optimizer score
  (low/mid/high) to evaluate difficulty-dependent performance.
\end{enumerate}

\textbf{Hard rule}: No parent-leaf-level leakage across train/val/test.

\subsection{Evaluation Metrics}

\paragraph{Parameter accuracy.}
Mean absolute error (MAE), root mean squared error (RMSE), and $R^2$
per parameter and aggregate. Parameters normalized to $[0,1]$ using bounds.

\paragraph{Physics score recovery.}
$\cS(p, \hat{\btheta}_p) / \cS(p, \btheta^*_p)$ compares predicted
to oracle scores using the actual physics scorer.

\paragraph{Runtime.}
Wall-clock time per patch, scorer evaluations per patch,
speedup ratio vs.\ full optimizer.

\paragraph{Calibration.}
Binned calibration (10 uncertainty bins) and risk-coverage curves.

\paragraph{Aggregated leaf-level accuracy.}
MAE of aggregated patch predictions vs.\ full-leaf oracle parameters.

%% ============================================================
\section{Results}
\label{sec:results}

\subsection{Baseline Comparison}

Table~\ref{tab:baselines} summarizes patch-level test performance across
all methods and backbones.

\begin{table}[h]
\centering
\caption{Patch-level test-set results (220 test patches from 44 test leaves).
Aggregate MAE across 7 normalized parameters.}
\label{tab:baselines}
\begin{tabular}{llcc}
\toprule
Method & Features & Test MAE & Test $R^2$ \\
\midrule
B0: Constant mean     & ---         & 1.68  & $0.00$ \\
B1: Ridge             & Handcrafted & 1.65  & $-0.05$ \\
B1: Ridge             & DINOv2      & 3.06  & $-2.91$ \\
B1: Ridge             & ResNet      & ---   & ---    \\
B2: MLP               & DINOv2      & 6.39  & $-1832.6$ \\
B3: Ensemble          & DINOv2      & 5.97  & $-183.3$ \\
\midrule
\multicolumn{4}{l}{\emph{External domain (18.7), ensemble}} \\
B3: Ensemble          & DINOv2      & 5.90  & $-93.3$ \\
\bottomrule
\end{tabular}
\end{table}

The constant-mean baseline (MAE = 1.68) is strong due to the relatively
narrow parameter distribution across patches. The handcrafted ridge
predictor (MAE = 1.65) marginally improves over the mean baseline.
The DINOv2 ridge (MAE = 3.06) performs worse in aggregate MAE
but provides richer feature representations for ensemble uncertainty estimation.

The ensemble (MAE = 5.97) has higher point-prediction error but serves
as the uncertainty source for selective routing. As in the companion
whole-leaf study, the ensemble's primary value is in uncertainty-based routing
rather than raw prediction accuracy.

\subsection{Real Refinement Results}

Table~\ref{tab:refinement} shows the runtime--quality tradeoff across
six refinement budgets, evaluated on 407 test patches with real PDE scoring.

\begin{table}[h]
\centering
\caption{Patch-level refinement results (407 test patches).
All refinement uses actual physics scorer calls.}
\label{tab:refinement}
\begin{tabular}{lcccc}
\toprule
Budget & Mean Runtime (s) & Scorer Calls & MAE & Score Recovery \\
\midrule
Direct   & 0.17 & 1.0  & 5.34 & 0.74 \\
Tiny     & 0.66 & 4.0  & 5.28 & 0.78 \\
Small    & 1.13 & 7.0  & 5.24 & 0.80 \\
Medium   & 1.76 & 11.0 & 5.22 & 0.81 \\
Large    & 3.34 & 21.0 & 5.21 & 0.82 \\
Full     & 2.84 & 23.1 & 2.11 & 1.01 \\
\bottomrule
\end{tabular}
\end{table}

Full refinement achieves dramatic MAE improvement (5.34 $\to$ 2.11) with
score recovery $>1.0$ (slightly exceeding the oracle on some patches).
Intermediate budgets (tiny--large) show modest MAE improvement but
consistent score recovery gains, confirming that even a few real PDE
evaluations provide useful signal.

Figure~\ref{fig:pareto} shows the runtime--quality Pareto frontier.

\begin{figure}[h]
\centering
\includegraphics[width=0.7\textwidth]{fig01_runtime_quality_pareto.png}
\caption{Patch refinement: runtime--quality Pareto frontier across budgets.}
\label{fig:pareto}
\end{figure}

\subsection{Selective Routing Analysis}

We evaluated 136 selective routing policies on the 220 test patches.
The best policy achieves:
\begin{itemize}[leftmargin=*,itemsep=1pt]
\item Thresholds: $\tau_\text{low} = 0.556$, $\tau_\text{high} = 0.581$
\item Fraction direct: 4.5\%, fraction refine: 5.0\%, fraction fallback: 90.5\%
\item Mean MAE: 2.33 (vs.\ 5.97 direct-only)
\item Mean runtime: 2.69\,s per patch
\item Speedup: $11.1{\times}$ vs.\ full optimization on all patches
\end{itemize}

The high fallback fraction reflects the challenging nature of patch-level
prediction: uncertainty is generally high, pushing most patches through
full re-optimization. Nevertheless, the $11.1{\times}$ speedup demonstrates
that selective routing remains effective at the patch level.

\subsection{Patch-to-Leaf Aggregation}

Table~\ref{tab:aggregation} compares six aggregation operators on 44 test leaves.

\begin{table}[h]
\centering
\caption{Patch-to-leaf aggregation results (44 test leaves).
Leaf MAE measures the error between aggregated patch predictions
and full-leaf oracle parameters $\btheta^*_\text{leaf}$.}
\label{tab:aggregation}
\begin{tabular}{lc}
\toprule
Aggregation Method & Mean Leaf MAE \\
\midrule
Selective mixture        & \textbf{7.99} \\
Unweighted mean          & 10.02 \\
Uncertainty-weighted     & 10.02 \\
Lesion-fraction-weighted & 10.02 \\
Top-$k$ confident ($k{=}3$) & 10.04 \\
Robust median            & 10.03 \\
\bottomrule
\end{tabular}
\end{table}

\begin{figure}[h]
\centering
\includegraphics[width=0.75\textwidth]{fig02_aggregation_comparison.png}
\caption{Patch-to-leaf aggregation method comparison (44 test leaves).}
\label{fig:aggregation}
\end{figure}

Selective mixture (leaf MAE = 7.99) significantly outperforms all non-routed
methods (${\sim}$10.02). This confirms that the routing-aware aggregation---using
direct predictions for confident patches, refined estimates for uncertain ones,
and oracle fallback for the most uncertain---transfers meaningful quality
gains from the patch to the leaf level.

The near-identical performance of all non-routed methods (mean, uncertainty-weighted,
lesion-weighted, top-$k$, median) suggests that, without routing information,
the aggregation operator choice matters less than the quality of individual
patch estimates.

\subsection{Calibration and Risk-Coverage}

Figure~\ref{fig:calibration} shows the calibration plot (10 uncertainty bins).
The correlation between predicted uncertainty and actual error is positive
but noisy, reflecting the challenge of calibrating uncertainty at the patch level.

\begin{figure}[h]
\centering
\includegraphics[width=0.48\textwidth]{fig04_calibration.png}
\hfill
\includegraphics[width=0.48\textwidth]{fig03_risk_coverage.png}
\caption{Left: Uncertainty calibration. Right: Risk-coverage tradeoff.
Analysis on 220 test patches.}
\label{fig:calibration}
\end{figure}

The risk-coverage curve (Figure~\ref{fig:risk_coverage}) shows that rejecting
the most uncertain 10\% of patches reduces MAE from 5.97 to 5.76,
with monotonic risk reduction as coverage decreases. This validates
that ensemble uncertainty provides a meaningful ranking for selective behavior.

\subsection{Domain Robustness}

\begin{figure}[h]
\centering
\includegraphics[width=0.7\textwidth]{fig08_domain_robustness.png}
\caption{Domain comparison: primary leaves vs.\ external 18.7 patches.}
\label{fig:domain}
\end{figure}

The external 18.7 domain shows comparable ensemble MAE (5.90 vs.\ 5.97 on
primary test), suggesting that the learned features generalize across
collection batches. This is encouraging for deployment, though we note
that parameter-level ground truth is available for both domains (patch oracle),
allowing direct comparison.

\subsection{Phenotype Geometry}

\begin{figure}[h]
\centering
\includegraphics[width=0.48\textwidth]{fig05_geometry_pca.png}
\hfill
\includegraphics[width=0.48\textwidth]{fig12_geometry_3d.png}
\caption{Left: PCA of patch parameter space. Right: 3D PCA scatter.
Oracle, predicted, and refined parameter distributions.}
\label{fig:geometry}
\end{figure}

The convex hull analysis reveals:
\begin{itemize}[leftmargin=*,itemsep=1pt]
\item Oracle hull volume: 15.59
\item Predicted hull volume: 2.73 (82\% contraction)
\item Refined hull volume: 0.76 (95\% contraction)
\end{itemize}

Centroid shifts: oracle $\to$ predicted = 39.60, oracle $\to$ refined = 36.12.
The substantial hull contraction indicates that the predictor concentrates
predictions near the mean, consistent with the strong constant-mean baseline.
Refinement further concentrates predictions (improving individual patch scores
but reducing distributional coverage).

\subsection{Within-Leaf Spatial Heterogeneity}
\label{sec:heterogeneity}

\begin{figure}[h]
\centering
\includegraphics[width=0.7\textwidth]{fig07_spatial_heterogeneity.png}
\caption{Distribution of within-leaf parameter heterogeneity across 174 leaves.}
\label{fig:heterogeneity}
\end{figure}

Across 174 leaves with $\geq 2$ oracle-optimized patches each,
the mean within-leaf parameter standard deviation is $\sigma_\text{intra} = 1.70$.
The parameter with the highest within-leaf variability is consistently
$E_\text{thr}$ (energy threshold), reflecting the spatial heterogeneity
of lesion severity within individual leaves.

This finding is a key contribution: it demonstrates that a single
parameter vector per leaf is an oversimplification. Different regions
of the same leaf have genuinely different biophysical states, and
patchwise inference captures this spatial structure.

\subsection{Per-Parameter Difficulty}

\begin{figure}[h]
\centering
\includegraphics[width=0.7\textwidth]{fig09_per_parameter_difficulty.png}
\caption{Per-parameter prediction difficulty (ensemble, test patches).}
\label{fig:per_param}
\end{figure}

The energy threshold ($E_\text{thr}$) is by far the most difficult parameter
to predict (MAE = 39.6), which is expected given its large dynamic range
$[5, 60]$ and high within-leaf variability.
The remaining parameters have MAE between 0.21 and 0.62.

\subsection{Qualitative Case Studies}

Figure~\ref{fig:qualitative} shows representative patch panels with
oracle, direct, medium-refined, and full-refined segmentation masks.

\begin{figure}[h]
\centering
\includegraphics[width=\textwidth]{qualitative_case_1_panel.png}
\includegraphics[width=\textwidth]{qualitative_external_panel.png}
\caption{Top: primary leaf patch case study. Bottom: external 18.7 patch.
Left to right: raw patch, oracle mask, direct prediction mask,
medium refinement mask, full refinement mask.}
\label{fig:qualitative}
\end{figure}

%% ============================================================
\section{Ablation Studies}
\label{sec:ablations}

\subsection{Backbone Comparison}

Table~\ref{tab:backbone} compares the three feature backbones.

\begin{table}[h]
\centering
\caption{Backbone comparison (test split, ridge regression).}
\label{tab:backbone}
\begin{tabular}{lcc}
\toprule
Backbone & Test MAE (Aggregate) & Test $R^2$ \\
\midrule
Handcrafted ($d{=}45$) & 1.65 & $-0.048$ \\
DINOv2 ($d{=}384$)     & 3.06 & $-2.91$ \\
ResNet-18 ($d{=}512$)  & --- & --- \\
\bottomrule
\end{tabular}
\end{table}

The handcrafted features outperform both deep backbones in aggregate MAE
for ridge regression, likely because the low-dimensional handcrafted
features ($d{=}45$) are less prone to overfitting with the limited
training sample (${\sim}520$ patches). However, the DINOv2 ensemble
provides the uncertainty estimates needed for selective routing.

\subsection{Refinement Budget}

\begin{figure}[h]
\centering
\includegraphics[width=\textwidth]{fig11_refinement_budget.png}
\caption{Refinement budget ablation: quality (left) and runtime (right).}
\label{fig:budget_ablation}
\end{figure}

The full budget achieves dramatically lower MAE (2.11) than intermediate
budgets (5.21--5.34). This large gap suggests that the search landscape
at the patch level requires substantial exploration to escape local optima,
and that intermediate budgets provide only marginal improvement over
direct prediction.

\subsection{Aggregation Method}

Selective mixture is the only aggregation method that substantially
outperforms the others (7.99 vs.\ ${\sim}$10.02), confirming that
routing information is the dominant factor in aggregation quality.

\subsection{Patch Extraction Parameters}

The current study uses a single patch configuration:
$224{\times}224$ patches, stride 112, $\leq 5$ patches per leaf (primary),
$\leq 3$ patches per leaf (external).
Future work should evaluate alternative patch sizes (e.g., $128$, $336$)
and stride ratios.

\subsection{Performance Comparison}

\begin{figure}[h]
\centering
\includegraphics[width=0.7\textwidth]{fig10_performance_comparison.png}
\caption{Aggregate test MAE across all methods.}
\label{fig:performance}
\end{figure}

\subsection{Uncertainty vs.\ Error}

\begin{figure}[h]
\centering
\includegraphics[width=0.65\textwidth]{fig06_uncertainty_vs_error.png}
\caption{Ensemble uncertainty vs.\ prediction error (test patches).}
\label{fig:unc_error}
\end{figure}

The positive correlation between uncertainty and error validates the
ensemble's utility as a routing signal, even though calibration is imperfect.

%% ============================================================
\section{Discussion}
\label{sec:discussion}

\subsection{The Patchification Argument}

The primary motivation for patchification is not merely to increase sample count
(from 174 leaves to 1{,}057 patches), but to \textbf{preserve fine-scale
lesion information} that is destroyed by downsampling.
The within-leaf heterogeneity analysis (Section~\ref{sec:heterogeneity})
provides direct evidence: with mean $\sigma_\text{intra} = 1.70$,
different patches from the same leaf genuinely require different
parameter configurations. A single whole-leaf parameter vector is an
averaging artifact that obscures this biophysical reality.

Patchification also enables compatibility with standard pretrained backbones
(DINOv2, ResNet-18) at their native input resolution, avoiding the
architectural modifications or multi-scale adaptations that would be needed
to process full-resolution images end-to-end.

\subsection{The Mean-Baseline Challenge}

As in the companion whole-leaf study, the constant-mean baseline is
remarkably competitive (MAE = 1.68), reflecting limited inter-patch
parameter variability relative to prediction noise.
This is a fundamental challenge for physics parameter prediction with
small sample sizes: the parameter distribution is narrow enough that
the mean is a strong estimator.

The ensemble predictor has higher MAE (5.97) but provides the uncertainty
estimates that enable selective routing. The selective policy achieves
$11.1{\times}$ speedup precisely because it identifies the minority of
patches where optimization matters.

\subsection{Aggregation: Route-Awareness Is Key}

The aggregation results (Table~\ref{tab:aggregation}) show that
\emph{how} patch estimates are obtained matters more than \emph{how}
they are combined. All non-routed methods perform nearly identically
(${\sim}$10.02 leaf MAE), while selective mixture (7.99) benefits from
having routed uncertain patches through refinement or oracle fallback.

This suggests that future work on patch-to-leaf aggregation should focus
on improving individual patch estimate quality (via better routing or
more efficient refinement) rather than on sophisticated aggregation operators.

\subsection{Refinement Landscape at Patch Scale}

The large gap between intermediate and full refinement (Table~\ref{tab:refinement})
suggests that the patch-level scoring landscape is challenging:
the reduced spatial extent of patches provides less signal per scorer
evaluation compared to whole-leaf images. This motivates research into
patch-specific refinement strategies (e.g., using multi-patch context
or learned refinement operators).

\subsection{Comparison with Companion Whole-Leaf Study}

The patchwise study mirrors the companion whole-leaf study in experimental
design but reveals complementary insights:
\begin{itemize}[leftmargin=*,itemsep=2pt]
\item \textbf{Unique to patches}: Within-leaf heterogeneity analysis,
  patch-to-leaf aggregation, spatial locality preservation.
\item \textbf{Shared methodology}: Three-tier selective routing,
  real short refinement, deep ensemble uncertainty, calibration/risk-coverage.
\item \textbf{Scale}: 1{,}057 patches vs.\ 174 leaves ($6{\times}$
  more data points for statistical analysis).
\item \textbf{Challenge}: Patch-level prediction is harder (larger MAE)
  because individual patches carry less global context than whole-leaf images.
\end{itemize}

\subsection{Limitations and Future Directions}

\paragraph{Patch size and stride.}
The current study uses a single patch configuration ($224{\times}224$, stride 112).
A systematic sweep over patch sizes and strides would characterize the
resolution--context tradeoff.

\paragraph{Sample size.}
With 174 primary leaves (870 patches), the effective training set remains
modest. Scaling to $400{+}$ leaves would improve model capacity and reduce
the dominance of the mean baseline.

\paragraph{Attention-based aggregation.}
The current aggregation operators are simple. Attention-based MIL
methods~\cite{ilse2018attention} could learn patch importance weights
end-to-end.

\paragraph{Conformal prediction.}
Conformal prediction~\cite{angelopoulos2023conformal} could provide
distribution-free coverage guarantees for patch-level parameter intervals.

\paragraph{Multi-objective routing.}
The current policy optimizes MAE subject to speedup constraints.
Joint optimization over multiple quality metrics (MAE, score recovery,
morphology fidelity) would provide richer deployment tradeoffs.

\paragraph{Overlapping patch consistency.}
Adjacent patches share spatial context. Enforcing consistency between
overlapping patch predictions (e.g., via spatial regularization)
could improve both patch-level and aggregated leaf-level accuracy.

%% ============================================================
\section{Conclusion}
\label{sec:conclusion}

We presented patchwise selective semi-amortized inference for physics-based
lesion phenotyping, extending the companion whole-leaf framework to operate
at the patch level. By extracting $224{\times}224$ patches from full-resolution
leaf images, we preserve fine-scale lesion boundaries, local texture,
and within-leaf heterogeneity that are destroyed by whole-leaf downsampling.

On 1{,}057 patches from 237 maize leaves (174 primary + 63 external),
with 7 physics parameters per patch, we demonstrated:
\begin{enumerate}[leftmargin=*,itemsep=2pt]
\item Real patch-level oracle optimization (870 + 187 patches),
  revealing substantial within-leaf parameter heterogeneity
  (mean $\sigma_\text{intra} = 1.70$).
\item Patch-level selective routing with $11.1{\times}$ speedup
  over full optimization.
\item Route-aware selective-mixture aggregation (leaf MAE = 7.99)
  outperforming all non-routed methods (${\sim}$10.02).
\item Complete calibration, risk-coverage, geometry, and
  domain robustness analyses.
\end{enumerate}

The patchwise approach is a natural companion to whole-leaf inference:
it trades computational cost for spatial resolution, enabling
local phenotype characterization that is impossible under global averaging.
Together, the two frameworks provide a complete toolkit for
physics-based lesion phenotyping at multiple spatial scales.

%% ============================================================
\section*{Acknowledgements}

Computational resources provided by [institution HPC cluster].

\bibliographystyle{plain}
\bibliography{../bib/references_patch_v2}

%% ============================================================
\appendix

\section{Additional Figures}
\label{app:figures}

\begin{figure}[h]
\centering
\includegraphics[width=\textwidth]{qualitative_case_2_panel.png}
\caption{Additional case study panel (primary leaf patch 2).}
\end{figure}

\begin{figure}[h]
\centering
\includegraphics[width=\textwidth]{qualitative_case_3_panel.png}
\caption{Additional case study panel (primary leaf patch 3).}
\end{figure}

\begin{figure}[h]
\centering
\includegraphics[width=\textwidth]{qualitative_case_4_panel.png}
\caption{Additional case study panel (primary leaf patch 4).}
\end{figure}

\section{Selective Policy Grid}
\label{app:policy_grid}

The full 136-policy grid search evaluated
$(q_\text{low}, q_\text{high})$ combinations across
$q_\text{low} \in \{0.05, 0.10, \ldots, 0.50\}$ and
$q_\text{high} \in \{0.15, 0.20, \ldots, 0.90\}$.
The best policy ($\tau_\text{low}=0.556$, $\tau_\text{high}=0.581$)
achieves a balance between direct acceptance (4.5\%),
refinement (5.0\%), and fallback (90.5\%), with MAE = 2.33
and $11.1{\times}$ speedup.

\section{Patch Extraction Details}
\label{app:extraction}

\begin{itemize}[leftmargin=*,itemsep=2pt]
\item Patch size: $224 \times 224$ pixels
\item Stride: 112 pixels (50\% overlap)
\item Content filter: $\geq$10\% non-background pixels
\item Primary leaves: $\leq 5$ patches per leaf, lesion-aware selection
\item External leaves: $\leq 3$ patches per leaf
\item Total: 870 primary + 187 external = 1{,}057 patches
\end{itemize}

\section{Within-Leaf Heterogeneity Details}
\label{app:heterogeneity}

The energy threshold ($E_\text{thr}$) dominates within-leaf variability
in all 174 multi-patch leaves. This is consistent with the
physical interpretation: lesion severity varies spatially within a leaf,
producing regions with very different energy landscapes that require
different threshold settings for accurate segmentation.

The mean within-leaf range (max $-$ min across patches) is:
\begin{itemize}[leftmargin=*,itemsep=1pt]
\item $E_\text{thr}$: range $\approx$ 10--17 (out of $[5, 60]$)
\item $\gamma$: range $\approx$ 0.1--0.3
\item $\mu, \lambda$: range $\approx$ 0.05--0.2
\item $\alpha$: range $\approx$ 0.01--0.03
\end{itemize}

\section{Full Backbone Comparison}
\label{app:backbone}

Three backbones were evaluated: handcrafted features ($d{=}45$),
DINOv2-ViT-S/14 with registers ($d{=}384$), and ResNet-18 ($d{=}512$).
For each backbone, Ridge and MLP regressors were trained.
The DINOv2 ensemble (5$\times$MLP) serves as the primary uncertainty source.

Handcrafted features include:
RGB channel means and standard deviations (6),
Lab channel statistics (6),
Gabor filter bank responses at 4 orientations and 3 scales (12),
Local Binary Pattern histogram (10),
edge density (Canny, 1),
Hu moments (7),
non-zero fraction (1),
mean gradient magnitude (1),
entropy (1).

\end{document}
